% =============================================================================
%  Supplementary Material
%  Lineage-Aware Memory Governance for Enterprise AI Agents
%  Accompanies the IEEE Access submission
% =============================================================================
\documentclass[9pt,twocolumn,twoside]{extarticle}

\usepackage[T1]{fontenc}
\usepackage[utf8]{inputenc}
\usepackage{mathptmx}
\usepackage[margin=0.75in,columnsep=0.25in,twoside]{geometry}
\usepackage{microtype}
\usepackage{amsmath,amssymb,amsthm}
\usepackage{mathtools}
\usepackage{graphicx}
\usepackage{booktabs}
\usepackage{array}
\usepackage{multirow}
\usepackage{xcolor}
\usepackage{tcolorbox}
\tcbuselibrary{listings,breakable,skins}
\usepackage{listings}
\usepackage{cite}
\usepackage[hidelinks]{hyperref}
\usepackage{cleveref}
\usepackage{enumitem}
\usepackage{caption}
\usepackage{fancyhdr}
\usepackage{pifont}
\usepackage{balance}
\usepackage{titlesec}
\usepackage{placeins}

\renewcommand{\dbltopfraction}{0.9}
\renewcommand{\dblfloatpagefraction}{0.85}
\setcounter{dbltopnumber}{4}
\renewcommand{\topfraction}{0.9}
\renewcommand{\textfraction}{0.07}
\renewcommand{\floatpagefraction}{0.8}

\captionsetup{font=small,labelfont=bf,labelsep=period}

\titleformat{\section}
  {\centering\normalfont\fontsize{10}{12}\bfseries}
  {\Roman{section}.}{0.6em}{\MakeUppercase}
\titleformat{\subsection}
  {\normalfont\bfseries\itshape}
  {\Alph{subsection}.}{0.5em}{}
\titlespacing*{\section}{0pt}{10pt plus 2pt minus 1pt}{6pt}
\titlespacing*{\subsection}{0pt}{7pt plus 1pt}{4pt}

\newtheorem{theorem}{Theorem}
\newtheorem{corollary}[theorem]{Corollary}
\newtheorem{definition}{Definition}
\newtheorem{assumption}{Assumption}
\newtheorem*{remark}{Remark}

\definecolor{algbg}{HTML}{F8FAFC}
\definecolor{algframe}{HTML}{1E3A5F}
\definecolor{algkw}{HTML}{1D4ED8}
\definecolor{algcmt}{HTML}{6B7280}
\definecolor{algnum}{HTML}{374151}

\lstdefinelanguage{pseudocode}{
  morekeywords={for,if,else,return,true,false,and,or,not,in,to,do,
                begin,end,each,all,break,continue,of,with,not\_subset},
  keywordstyle=\bfseries\color{algkw},
  morecomment=[l]{\%\ },
  commentstyle=\itshape\color{algcmt},
  basicstyle=\ttfamily\scriptsize,
  numbers=left,
  numberstyle=\tiny\color{algnum},
  numbersep=6pt,
  xleftmargin=14pt,
  framexleftmargin=14pt,
  breaklines=true,
  showstringspaces=false,
  tabsize=2,
}

\newtcblisting[auto counter]{algorithmbox}[3][]{%
  title={\textbf{Algorithm~\thetcbcounter}: #2\hfill\textit{\small#3}},
  colback=algbg, colframe=algframe, fonttitle=\small\bfseries,
  listing only, listing options={language=pseudocode},
  breakable, enhanced, #1
}

\definecolor{myblue}{HTML}{1D4ED8}

\pagestyle{fancy}
\fancyhf{}
\fancyhead[LE]{\small\itshape Supplementary Material --- Sangaraju and Vissa}
\fancyhead[RO]{\small\itshape Supplementary Material --- Sangaraju and Vissa}
\fancyfoot[C]{\small\thepage}
\renewcommand{\headrulewidth}{0.4pt}
\renewcommand{\footrulewidth}{0pt}

\fancypagestyle{firstpage}{%
  \fancyhf{}
  \fancyfoot[C]{\small\thepage}
  \renewcommand{\headrulewidth}{0pt}
  \renewcommand{\footrulewidth}{0pt}
}

\setlength{\parindent}{1em}

\begin{document}

\twocolumn[
\begin{center}
{\LARGE\bfseries Supplementary Material\par}
\vspace{6pt}
{\large Lineage-Aware Memory Governance: A Derivation-Gated Framework
for Privacy-Preserving Column-Level Access Control in Enterprise AI Agents\par}
\vspace{8pt}
{\large Venkata Sangaraju\textsuperscript{1} and Sudhir Vissa\textsuperscript{2}\par}
\vspace{5pt}
{\small
\textsuperscript{1}Independent Researcher
(ORCID: \href{https://orcid.org/0009-0001-7716-1342}{0009-0001-7716-1342})\\
\textsuperscript{2}SAGE7 AI, Georgetown, Texas, USA
(ORCID: \href{https://orcid.org/0009-0003-7865-0863}{0009-0003-7865-0863})\par}
\end{center}
\vspace{8pt}
\noindent\rule{\textwidth}{0.6pt}
\vspace{6pt}
\noindent
This document accompanies the IEEE Access submission
\emph{``Lineage-Aware Memory Governance: A Derivation-Gated Framework for
Privacy-Preserving Column-Level Access Control in Enterprise AI Agents.''}
It reproduces, in full, material that the main paper summarises for length:
extended related-work discussion (\S\ref{sec:supp-related}), the
line-by-line walkthrough of Algorithm~2 (\S\ref{sec:supp-algo}; Algorithm~1's
walkthrough and the full proof of Theorem~1 are now given directly in the
main paper and are only reproduced here for self-containedness), extended
threat descriptions (\S\ref{sec:supp-threat}), the unevaluated composite
conflict-detector design (\S\ref{sec:supp-d4}), and complete per-experiment
result tables (\S\ref{sec:supp-tables}).
All numbering in this document is local to this document; cross-references
to the main paper are given as plain section numbers (main paper, \S$n$).
\vspace{6pt}
\noindent\rule{\textwidth}{0.6pt}
\vspace{10pt}
]

\thispagestyle{firstpage}

% =============================================================================
\section{Extended Related Work}
\label{sec:supp-related}
% =============================================================================
The main paper (\S2) summarises related work in one page against six
governance dimensions (Table 1). This section reproduces the full
per-system discussion.

\subsection{Agent Memory Systems}

Du~\cite{du2026survey} formalises agent memory as a write-manage-read loop
with a three-dimensional taxonomy spanning temporal scope, representational
substrate, and control policy.
MemGPT~\cite{packer2023memgpt} instantiated the hierarchy by treating the LLM
as an OS that pages between a bounded context window and an external store via
function calls, demonstrating coherent long-horizon conversations and
large-document analysis.
\emph{Limitation relative to our work:} governance in MemGPT is implicit---the
LLM decides what to write and retrieve---with no lineage tracking at any layer.

Zep~\cite{rasmussen2025zep} introduced a temporal knowledge graph with dual-axis
timestamping (event time vs.\ ingestion time), enabling point-in-time queries.
\emph{Limitation:} its graph captures \emph{when} information arrived, not
\emph{how} it was derived.
A-MEM~\cite{xu2025amem}---a NeurIPS 2025 paper---builds Zettelkasten-style
linked notes with dynamic indexing, achieving strong benchmark performance across
six foundation models.
\emph{Limitation:} linking structure encodes semantic relatedness between notes,
not the data-derivation path of a computed result.
All three systems represent the state of the art in conversational or
task-level memory and are well-suited to remembering facts, past
conversations, and user preferences, but none model the provenance of
computed analytical results.

\subsection{Governed and Enterprise Memory Architectures}

Taheri's Governed Memory~\cite{taheri2026governed} addresses a real enterprise
gap---dozens of agents operating with no shared memory and no common
governance---through a dual memory model, tiered governance routing, and a
closed-loop schema lifecycle, achieving 99.6\% fact recall and zero
cross-entity leakage across 500 adversarial queries in a 250-event controlled
experiment.
\emph{Limitation:} its governance operates on \emph{what is stored and who may
access it as a whole entry}; it does not track derivation, so it cannot detect
that a permitted entry embeds a forbidden column.
Lam et al.'s SSGM~\cite{lam2026ssgm} decouples memory evolution from execution
by enforcing consistency verification, temporal decay modelling, and dynamic
access control before any memory consolidation.
\emph{Limitation:} SSGM's dynamic access control enforces agent-topology
constraints (who may communicate with whom), not column-level derivation
constraints (what data was touched to produce a result)---the two mechanisms
are complementary rather than overlapping.
Oracle AI Agent Memory~\cite{oracle2026memory} provides a commercially
deployed governed memory core with fine-grained lifecycle management and
role-based access controls.
\emph{Limitation:} role assignment controls access to a memory entry as a
whole, not to the individual columns that contributed to it.
Abaev et al.'s AgentGuardian~\cite{abaev2026agentguardian} learns
context-aware access-control policies from execution traces to govern which
tool calls an agent may issue.
\emph{Limitation:} it governs \emph{tool-call} legitimacy at execution time,
not the provenance of values already committed to a shared memory store.

\subsection{Column-Level Security and Data Provenance}

Apache Ranger, Databricks Unity Catalog, and BigQuery column-level security
enforce column-gated access at the \emph{query layer}: they prevent an
unauthorised agent from issuing a query that touches a restricted column.
Subramaniam and Krishnan's DePLOI~\cite{subramaniam2024deploi} uses NL2SQL
task decomposition to synthesise and audit database access-control
implementations against natural-language intent, improving the
trustworthiness of the source-layer policies themselves.
\emph{Limitation relative to our work:} these are all source-layer controls
that fire at query time; none address a result that was legitimately computed
once, cached, and then served to a different, less-privileged requester.
The W3C PROV standard~\cite{moreau2011prov} encodes entity-activity-agent
relationships for data provenance, and Herschel et al.~\cite{herschel2017survey}
survey systems that use provenance to track data flow through complex
pipelines.
Column-level lineage extraction systems such as
LINEAGEX~\cite{zhang2025lineagex}---which parses SQL query trees to infer
column-level derivation with high coverage and accuracy, demonstrated at ICDE
2025---represent the state of the art in automatic lineage extraction.
Zhou et al.'s survey of the LLM$\times$DATA intersection~\cite{zhou2025llmxdata}
catalogues how data systems and LLMs increasingly interact but does not treat
agent memory retrieval as a governance surface in its own right.
\emph{Limitation:} the provenance literature focuses on raw database records
and ETL pipelines; our AMU schema applies the same column-level-derivation
concept to \emph{agent memory entries}, and LINEAGEX or
\textsc{SQLGlot}~\cite{sqlglot2023} (used in Experiment~6) are the natural
tools for automatically populating AMU lineage fields in production.
No prior work applies column-level provenance to agent memory retrieval
governance, or uses derivation graphs to gate cross-department memory
sharing---this is the precise gap we fill.

\subsection{Privacy and Security in Multi-Agent Systems}

El Yagoubi et al.'s AgentLeak~\cite{elyagoubi2026agentleak}---the first
full-stack benchmark for privacy leakage in multi-agent LLM systems---shows
that 68.8\% of leakage occurs through inter-agent channels (messages, shared
memory, tool arguments) invisible to output-only auditors, across 4,979
validated execution traces.
Naik et al.'s OMNI-LEAK~\cite{naik2026omnileak} demonstrates that a single
indirect prompt injection can compromise several agents and leak sensitive
data even with data-access controls in place; Alizadeh et
al.~\cite{alizadeh2025leak} show that simple prompt-injection payloads can
cause an LLM agent to leak personal data it observed during task execution,
and Wang et al.'s taxonomy of prompt-injection
threats~\cite{wang2026landscape} and Dehghantanha and Homayoun's SoK of the
agentic attack surface~\cite{dehghantanha2026sok} both catalogue memory and
context poisoning as first-class attack vectors against agents.
Nguyen et al.'s survey of multi-agent system security~\cite{nguyen2026security}
derives a 193-item threat taxonomy from production multi-agent architectures,
including threats to shared memory and inter-agent trust.
Summers and Wu's Data Flow Control~\cite{summers2026dataflow} enforces
declarative, tuple-level data-safety policies inside the DBMS query layer
itself using provenance monomials---closely related in spirit to our
column-level gate, but operating on live query execution rather than on a
persistent cross-agent memory cache.
Wang et al.'s AuthGraph~\cite{wang2026authgraph} defends individual agents
against prompt-injection by structurally comparing an execution-derived
provenance graph against an authorisation graph, reducing attack success rates
from $\sim$40\% to 1--2\% on two benchmarks---the closest prior work to our
own in combining a provenance graph with an authorisation check, but scoped to
single-agent tool-calling safety rather than cross-department memory
retrieval.
Naseh et al.~\cite{naseh2025riddle} show that retrieval-augmented generation
systems are vulnerable to stealthy membership-inference attacks that recover
whether a specific record was present in a retrieval corpus, a threat that
targets what is stored rather than how it was derived.
\emph{Limitation relative to our work:} this literature identifies leakage
through inter-agent communication channels, prompt-level manipulation of a
single agent, and inference over stored content; our work identifies a
complementary vector---leakage through a \emph{shared analytical memory
store} that persists even when every individual communication link and every
individual agent is otherwise correctly secured.
The two families of failure modes are independent and both need to be
addressed in production; the threat model (main paper, \S5) (threat T8) discusses how
prompt-injection and inference-style attacks interact with the lineage record
itself.

\subsection{Metric-Definition Governance}

Semantic drift---divergent definitions of the same KPI across
teams---is a documented source of dashboard inconsistency in enterprise
BI~\cite{colrows2024metrics}.
Semantic-layer tools (dbt Metrics, Databricks Unity Catalog, Looker LookML)
address this for human analysts by centralising metric definitions, and the
Open Semantic Interchange initiative~\cite{snowflake2025osi}, launched in
September 2025 by Snowflake, Salesforce, dbt Labs, BlackRock, and
RelationalAI, is an industry-wide push toward a vendor-neutral metric
specification.
\emph{Limitation relative to our work:} these tools govern metrics that
\emph{humans} define centrally in advance; they provide no mechanism for
detecting drift when an autonomous agent computes a KPI on the fly through
its own generated query.
Our \texttt{definition\_hash} conflict detector (Algorithm~\ref{alg:write})
provides a lightweight, retrofittable foundation for automated metric
governance in agentic systems that write results outside a pre-registered
semantic layer; the extended fuzzy experiment in the extended fuzzy conflict-detection experiment (main paper, \S6)
characterises where hash-based detection is, and is not, the right
granularity.

% =============================================================================
\section{Algorithm Walkthroughs}
\label{sec:supp-algo}
% =============================================================================
The main paper now gives the full line-by-line walkthrough for Algorithm~1
(Lineage-Gated Retrieve) directly in System Design (main paper, \S3.2),
including its worst-/best-case complexity derivation, so it is not repeated
here; the pseudocode is reproduced below only so that line numbers cited by
the proof in \S\ref{sec:supp-proof} remain unambiguous in this standalone
document. Algorithm~2 (Lineage-Aware Write with Conflict Detection) is
reproduced below with its full walkthrough, since the main paper gives only
a summary of it.

% ── Algorithm 1 (pseudocode only; walkthrough is in the main paper) ─────────
\begin{algorithmbox}[label={alg:retrieve}]{Lineage-Gated Retrieve}{%
  \textnormal{Input: \textit{store}, \textit{metric}, \textit{dept}, \textit{fresh}}%
  \quad Output: \textit{result}; \textit{leaked}=\textbf{False} guaranteed}
candidates <- store[metric]         % all stored AMUs for this metric
permitted  <- P(dept)               % column set the department may see
blocked    <- False

for each a in REVERSE(candidates):  % most-recent-first
  if S(a) not_subset permitted:     % derivation touches a forbidden column
    blocked <- True
    continue                        % skip this candidate; try the next
  end if
  return (reused=True, value=a.value, leaked=False, blocked=False)
end for

store[metric].append(fresh)
return (reused=False, value=fresh.value, leaked=False, blocked=blocked)
\end{algorithmbox}

% ── Algorithm 2 ──────────────────────────────────────────────────────────────
\begin{algorithmbox}[label={alg:write}]{Lineage-Aware Write with Conflict Detection}{%
  \textnormal{Input: \textit{store}, AMU \textit{a}} \quad
  Output: \textit{conflict} \textbf{bool}}
% store    : shared memory dict (mutated in place)
% a        : new AMU to persist
% conflict : True iff a cross-department definition conflict is detected

existing <- store[a.metric]   % all AMUs stored for this metric name
conflict <- False

for each e in existing:
  if e.dept != a.dept:        % cross-department pair
    if H(e) != H(a):          % definition hashes differ -> potential conflict
      conflict <- True
      break                   % one conflict is sufficient; raise alert
    end if
  end if
end for

store[a.metric].append(a)     % always persist; conflict is surfaced, not blocked
return conflict
\end{algorithmbox}

\paragraph{Line-by-line walkthrough.}
Line~5 gathers every AMU already stored under the same metric name as the
incoming write, regardless of which department produced it.
Line~9 restricts the comparison to \emph{cross-department} pairs: two AMUs
written by the same department under the same metric name are assumed to
share a definition by construction and are not compared.
Line~10 compares the cached definition hashes $\mathrm{H}(\cdot)$
(Equation~\ref{eq:defhash}) rather than the raw lineage graphs, which is what
makes the comparison $O(1)$ per pair rather than requiring a graph-isomorphism
check.
Lines~11--12 raise the conflict flag and exit the loop on the first
cross-department hash mismatch found; one divergent definition is sufficient
evidence of drift, so the algorithm does not need to enumerate all of them.
Line~17 always persists the new AMU, whether or not a conflict was detected:
\emph{the write path never blocks on a conflict}, it only surfaces one for
downstream alerting (the threat model (main paper, \S5), T2).
This design choice trades a small amount of governance strictness for
operational safety---a metric-definition disagreement should never cause a
legitimate business query to fail outright.
\emph{Complexity:} the loop performs at most $n$ hash comparisons, each
$O(1)$, giving $O(n)$ worst case (no conflict, full scan) and $O(1)$ best case
(conflict found on the first cross-department comparison), matching the
empirical measurements in Table~\ref{tab:runtime}.

% =============================================================================
\section{Full Proof of the Sensitive-Column Safety Theorem}
\label{sec:supp-proof}
% =============================================================================
The main paper (\S4) now states Theorem 1, its full proof, and its two
corollaries in full (not merely a proof sketch). The complete statements and
proof are reproduced below as well, so that this document remains readable
independently of the main paper.
Equations (1)--(2) define the cached sensitivity tag $\mathrm{S}(a)$ and
definition hash $\mathrm{H}(a)$ referenced throughout.

\begin{align}
  \mathrm{S}(a) &= \bigcup_{i=1}^{k} C_i \;\cap\; \mathcal{S}
  \label{eq:sensitivity} \\[2pt]
  \mathrm{H}(a) &= \mathrm{SHA256}\!\left(
    \mathrm{sort}(\mathit{tables}) \,\|\,
    \mathrm{sort}\!\left(\bigcup_i C_i\right) \,\|\, \varphi
  \right)_{\![0{:}12]}
  \label{eq:defhash}
\end{align}

\begin{assumption}[Complete Lineage]
\label{ass:complete}
For every AMU $a$ in the store, $a.\mathit{lineage}$ records every table, every
column, and the filter logic of the computation that produced $a.\mathit{value}$.
\end{assumption}

\begin{theorem}[Sensitive-Column Safety]
\label{thm:safety}
Under Assumption~\ref{ass:complete}, Algorithm~\ref{alg:retrieve} returns no
result to department $d$ that was derived from a sensitive column
$c \in \mathcal{S}$ with $c \notin \mathrm{P}(d)$.
\end{theorem}

\begin{remark}[Scope of the guarantee]
\label{rem:scope}
The gate checks $\mathrm{S}(a) \subseteq \mathrm{P}(d)$ where
$\mathrm{S}(a) = \bigl(\bigcup_i C_i\bigr) \cap \mathcal{S}$ contains only
columns in the sensitive registry $\mathcal{S}$.
Non-sensitive columns that appear in a department's blocked list (such as
\texttt{amount} for Support in Table~\ref{tab:perms}) are \emph{not}
gate-enforced by Algorithm~\ref{alg:retrieve} as written.
An operator who wishes to gate-enforce non-sensitive columns as well should
extend $\mathcal{S}$ to include them, or change the gate condition to
$\bigl(\bigcup_i C_i\bigr) \subseteq \mathrm{P}(d)$ at the cost of
blocking more memory reuse.
\end{remark}

\begin{proof}
This proof is also given, in identical form, in the main paper (\S4); it is
reproduced here so the Supplementary Material remains self-contained.
By inspection, Algorithm~\ref{alg:retrieve} has exactly two statements that
return a value to the caller: line~10 (memory-served) and line~14 (fallback).
Every invocation terminates by reaching one of the two, since the
\texttt{for} loop either returns early at line~10 or exhausts \textit{candidates}
and falls through to lines~13--14.
We show the returned AMU's derivation touches no forbidden sensitive column
on each path.

\smallskip\noindent\textbf{Case 1: memory-served path (line~10).}
Reaching line~10 for candidate $a$ requires the guard at line~6 to have
evaluated false, i.e.\ $\mathrm{S}(a) \subseteq \mathrm{P}(d)$ holds for that
$a$.
Unfolding Equation~\ref{eq:sensitivity},
$\mathrm{S}(a) = \bigl(\bigcup_{i=1}^{k} C_i\bigr) \cap \mathcal{S}$, so
$\mathrm{S}(a) \subseteq \mathrm{P}(d)$ is equivalent to: for every step
$s_i \in a.\mathit{lineage}$ and every column $c \in C_i \cap \mathcal{S}$,
$c \in \mathrm{P}(d)$.
Contrapositively, there is no sensitive column $c \in \mathcal{S}$ with
$c \notin \mathrm{P}(d)$ appearing in any $C_i$---i.e., no such $c$ appears
anywhere in $a$'s derivation.
Since Assumption~\ref{ass:complete} guarantees $a.\mathit{lineage}$ records
\emph{every} column touched by the computation that produced $a.value$, this
is precisely the safety property claimed for the value returned at line~10.

\smallskip\noindent\textbf{Case 2: fallback path (line~14).}
The input precondition to Algorithm~\ref{alg:retrieve} (stated in the
signature) requires that \textit{fresh} is computed entirely within
$\mathrm{P}(d)$ by the caller, before Algorithm~\ref{alg:retrieve} is
invoked.
Under Assumption~\ref{ass:complete}, \textit{fresh}.\textit{lineage} records
every column of that computation, so every column recorded is, by the
precondition, in $\mathrm{P}(d)$; in particular no sensitive column outside
$\mathrm{P}(d)$ is touched.

\smallskip\noindent
Both cases establish that the AMU returned to department $d$ has a derivation
containing no sensitive column outside $\mathrm{P}(d)$, for an arbitrary
invocation with arbitrary \textit{store} contents.
Since the two cases are exhaustive over all return statements, the claim
holds for every invocation of Algorithm~\ref{alg:retrieve}.
\end{proof}

\begin{corollary}[Robustness to unsafe store contents]
\label{cor:unsafe-contents}
The guarantee of Theorem~\ref{thm:safety} holds even when \textit{store}
contains AMUs whose derivation paths touch sensitive columns
$c \in \mathcal{S}$ with $c \notin \mathrm{P}(d)$: such candidates fail the
guard at line~6 for department $d$'s request and are skipped, never
returned, regardless of how many such unsafe candidates exist or where they
occur in \textit{candidates}.
This matters operationally because a shared store necessarily contains AMUs
written by \emph{every} department, including ones whose lineage is
unsafe for the department currently requesting: the guarantee does not
depend on the store being curated per department.
\end{corollary}

\begin{corollary}[Behaviour under policy updates]
\label{cor:policy-update}
Let $\mathrm{P}_1(d)$ and $\mathrm{P}_2(d)$ be two successive versions of
department $d$'s permission set, and let $\mathcal{S}_1, \mathcal{S}_2$ be
two successive versions of the sensitive-column registry.
If $\mathrm{S}(a)$ is cached using $\mathcal{S}_1$ at write time and the
organisation transitions to $(\mathcal{S}_2, \mathrm{P}_2)$ without
recomputing $\mathrm{S}(a)$ for existing AMUs, Theorem~\ref{thm:safety}'s
guarantee holds only with respect to $(\mathcal{S}_1, \mathrm{P}_1)$, not
$(\mathcal{S}_2, \mathrm{P}_2)$: a column newly added to $\mathcal{S}_2$ that
appears in $a$'s lineage but was not in $\mathcal{S}_1$ will not appear in the
stale cached $\mathrm{S}(a)$, and a candidate that should now be blocked may
instead pass the gate.
The guarantee is restored for all subsequent retrievals as soon as
$\mathrm{S}(a)$ is recomputed against $\mathcal{S}_2$ for every AMU with a
write epoch preceding the policy change---an $O(N)$ one-time cost over the
$N$ AMUs in the store, after which each retrieval again costs $O(n)$ as in
Theorem~\ref{thm:safety}.
In practice we recommend the lazy recomputation strategy of
the implementation considerations (main paper, \S3): recompute $\mathrm{S}(a)$ the first time an AMU
is touched after a policy-version bump, amortising the $O(N)$ cost over
subsequent retrievals rather than paying it eagerly.
This corollary formalises threat T3 (the threat model (main paper, \S5)).
\end{corollary}

\begin{remark}[Assumption dependency]
Assumption~\ref{ass:complete} is critical.
If lineage is reported by the computing agent rather than extracted automatically
from query execution plans, a buggy or adversarial agent can omit a column,
breaking the guarantee.
Automatic extraction via systems such as LINEAGEX~\cite{zhang2025lineagex} is
the appropriate production countermeasure; this dependency is also catalogued
as threat T5 in the threat model (main paper, \S5).
The completeness-degradation experiment (main paper, \S6) empirically quantifies how the guarantee degrades
as a function of lineage completeness.
\end{remark}

% =============================================================================
\section{Extended Threat Descriptions}
\label{sec:supp-threat}
% =============================================================================
The main paper (\S5, Table 4) summarises all eight threats in tabular form.
The full prose description of each threat, reproduced from the original
threat-model analysis, follows.

\subsection{Threats Addressed}
\subsection{Threats Addressed}

\begin{description}[leftmargin=1.2em,itemsep=3pt]
\item[T1 --- Sensitive-column leakage via cached result.]
  An agent retrieves a result derived from a sensitive column
  $c \in \mathcal{S}$ it is not permitted to see.
  \textit{Mitigation:} Algorithm~\ref{alg:retrieve} blocks such retrievals
  via the $\mathrm{S}(a) \subseteq \mathrm{P}(d)$ gate.
  Non-sensitive columns outside $\mathrm{P}(d)$ are not gate-enforced by
  the current design; see Remark~\ref{rem:scope}.

\item[T2 --- Silent metric definition drift.]
  Two teams store conflicting definitions of the same KPI under the same name.
  \textit{Mitigation:} Algorithm~\ref{alg:write} surfaces the divergence on
  write via the definition-hash conflict detector.

\item[T3 --- Stale sensitive data reuse.]
  An AMU derived from data that was sensitive at write time is retrieved after
  column permissions are updated.
  \textit{Mitigation:} the \texttt{epoch} field enables TTL-based eviction;
  sensitivity tags should be re-evaluated on retrieval when policies change.
\end{description}

\subsection{Threats Not Addressed}
\subsection{Threats Not Addressed}

\begin{description}[leftmargin=1.2em,itemsep=3pt]
\item[T4 --- Inference attacks.]
  A derived feature that \emph{encodes} sensitive information without naming
  the source column (e.g., a risk score computed from \texttt{income}) will not
  be caught unless the dependency is explicitly recorded in the lineage graph.
  This is the most important production gap.

\item[T5 --- Malicious lineage fabrication.]
  A rogue agent writes an AMU with an intentionally incomplete lineage graph,
  bypassing the gate.
  The degradation experiment (the completeness-degradation experiment (main paper, \S6)) quantifies how even
  \emph{accidental} under-reporting erodes the guarantee.
  Automatic lineage extraction is the appropriate countermeasure.

\item[T6 --- Aggregate inference.]
  A sequence of individually safe retrievals may allow reconstruction of
  sensitive values.
  Algorithm~\ref{alg:retrieve} is a per-retrieval check; it provides no
  differential privacy guarantee.

\item[T7 --- Insider threats.]
  An operator with direct store access can read or modify entries.
  The mechanism assumes a trusted store layer.

\item[T8 --- Prompt-injection and timing attacks on the lineage record.]
  Two related attack surfaces are specific to the lineage layer itself and are
  not covered by Theorem~\ref{thm:safety}.
  First, an indirect prompt injection embedded in tool output or retrieved
  content could manipulate an LLM-orchestrated agent into issuing a query that
  omits a sensitive join from its \emph{self-reported} lineage, or into
  requesting that a differently derived (but identically named) result be
  written under a misleading \texttt{metric\_name}; this is a lineage-specific
  instance of the broader agent attack surface catalogued
  by~\cite{wang2026landscape,dehghantanha2026sok,alizadeh2025leak,naik2026omnileak}.
  Second, differential timing of the retrieve path itself is a theoretical
  side channel: because the worst-case branch (all candidates blocked, $O(n)$)
  and the best-case branch (first candidate safe, $O(1)$) have measurably
  different latencies (Table~\ref{tab:runtime}), a co-located observer could in
  principle infer whether \emph{some} cross-department AMU exists for a given
  metric name, even without learning its value---an inference analogous to the
  membership-inference attacks demonstrated against retrieval-augmented
  generation systems by Naseh et al.~\cite{naseh2025riddle}.
  \textit{Mitigation:} for the prompt-injection surface, lineage should be
  extracted automatically from the executed query (as in Experiment~6,
  the real-agent experiment (main paper, \S6)) rather than self-reported by the agent, and defences such
  as AuthGraph's provenance/authorisation cross-check~\cite{wang2026authgraph}
  are complementary and can run upstream of the AMU write path.
  For the timing surface, constant-time gate evaluation (padding every
  retrieval to worst-case latency) closes the channel at a fixed
  $13.8\,\mu$s-per-retrieval cost, which Table~\ref{tab:runtime} shows is
  negligible relative to any real analytics query; we do not implement this
  hardening in the current prototype and flag it as required only for
  deployments with an explicit timing-side-channel threat in scope.
\end{description}

% =============================================================================
\section{Proposed Composite Conflict Detector (D4, Unevaluated)}
\label{sec:supp-d4}
% =============================================================================
The main paper's extended fuzzy conflict-detection experiment (\S6) reports
that the column-graph detector D3 is blind to logic-operator changes (the
``LO blind spot''). Motivated by this finding, we sketch a composite
detector D4 that combines D3's column-graph comparison with an explicit
structural check on filter logic. \textbf{D4 has not been implemented or
evaluated}; it is a design proposal reported as future work, precisely so
it can be benchmarked against the same 43-pair dataset before any
production claim is made about it. The design intent is that D4 subsumes
D3 exactly on True-Conflict and Threshold-Variant pairs (where column-graph
comparison alone is already correct) while adding sensitivity to
Logic-Operator pairs by comparing the \emph{operator multiset} of the
parsed filter expression rather than its literal token similarity.

\begin{algorithmbox}[label={alg:composite}]{Proposed Composite Conflict Detector (D4, unevaluated)}{%
  \textnormal{Input: lineage pair $(L_1, L_2)$} \quad Output: \textit{conflict} \textbf{bool}}
% L1, L2 : two lineage graphs for the same metric_name from different depts
% Combines D3 (column-graph) with an explicit filter-logic structural check.
% NOT implemented or evaluated in this paper -- design proposal only.

tables1, cols1 <- tables_and_columns(L1)
tables2, cols2 <- tables_and_columns(L2)

if tables1 != tables2 or cols1 != cols2:
  return True                        % D3's existing check: structural join/column diff

% Column-graphs are identical; D3 alone would return False here.
% D4 adds an explicit comparison of filter-logic structure.
ast1 <- parse_filter_logic(L1.varphi)  % e.g. via a boolean-expression parser
ast2 <- parse_filter_logic(L2.varphi)

if operator_multiset(ast1) != operator_multiset(ast2):
  return True                        % AND/OR/NOT composition differs: logic-operator conflict

if threshold_literals(ast1) != threshold_literals(ast2):
  return False                       % same logical structure, only a numeric threshold differs
                                      % (TV case: not a conflict, matches D3's original design intent)

return False                         % identical column-graph and identical filter-logic structure
\end{algorithmbox}

% =============================================================================
\section{Complete Per-Experiment Result Tables}
\label{sec:supp-tables}
% =============================================================================
The main paper (\S6) reports condensed results tables. The complete,
un-condensed tables for every experiment are reproduced below.

\subsection{Experiment 1: Department Permissions and Full Results}
\begin{table*}[h]
\centering
\caption{Department permissions (synthetic schema).
The gate (Algorithm~\ref{alg:retrieve}) enforces restrictions only for
columns in the sensitive registry
$\mathcal{S} = \{\texttt{ssn},\texttt{email},\texttt{income}\}$;
non-sensitive blocked columns (\texttt{amount}, \texttt{spend}) represent
the full policy intent but are not currently gate-enforced
(see Remark~\ref{rem:scope}).}
\label{tab:perms}
\small
\begin{tabular}{@{}p{0.14\linewidth}p{0.28\linewidth}p{0.46\linewidth}@{}}
\toprule
\textbf{Department} & \textbf{Permitted (summary)} & \textbf{Blocked (gate-enforced in $\mathcal{S}$)} \\
\midrule
Finance   & All columns            & None \\
Marketing & Commercial, order data & \texttt{ssn, income, email} \\
Support   & Customer, ticket meta  & \texttt{ssn, income, email} (\texttt{amount, spend}: policy only) \\
\bottomrule
\end{tabular}
\end{table*}

\begin{table*}[h]
\centering
\caption{Synthetic schema results (30 seeds, mean $\pm$ population std dev).
Bootstrap 95\% CI on naive$-$LA difference excludes zero for all comparisons
(Mann-Whitney $U$, $p < 0.001$).}
\label{tab:results1}
\small
\begin{tabular}{@{}p{0.24\linewidth}p{0.13\linewidth}p{0.15\linewidth}p{0.17\linewidth}p{0.16\linewidth}@{}}
\toprule
\textbf{Metric} & \textbf{No Memory} & \textbf{Naive Shared} & \textbf{Lineage-Aware} & \textbf{Diff CI (95\%)} \\
\midrule
Leak rate (\%)           & $0.0 \pm 0.0$ & $18.8 \pm 3.7$ & $0.0 \pm 0.0$\textsuperscript{$\dagger$} & $[+14.7, +22.8]$ \\
Memory reuse rate (\%)   & $0.0 \pm 0.0$ & $95.8 \pm 0.0$ & $82.6 \pm 2.3$                           & $[+11.2, +15.1]$ \\
Conflict recall (\%)     & n/a           & $0.0 \pm 0.0$  & $100.0 \pm 0.0$                          & --- \\
Blocked retrievals (seed=42) & 0         & 0              & 15                                       & --- \\
\bottomrule
\multicolumn{5}{l}{\textsuperscript{$\dagger$}Formal guarantee (Theorem~\ref{thm:safety}), not an empirical discovery.}
\end{tabular}
\end{table*}

\subsection{Experiment 2: TPC-H Full Results}
\begin{table*}[h]
\centering
\caption{TPC-H-inspired schema results (30 seeds, mean $\pm$ population std dev).
Both comparisons: $p < 0.001$, bootstrap 95\% CI excludes zero.}
\label{tab:results2}
\begin{tabular}{@{}lccc@{}}
\toprule
\textbf{Metric} & \textbf{Naive Shared} & \textbf{Lineage-Aware} & \textbf{Diff CI (95\%)} \\
\midrule
Leak rate (\%)            & $25.5 \pm 4.5$ & $0.0 \pm 0.0$\textsuperscript{$\dagger$} & $[+19.7, +31.2]$ \\
Memory reuse rate (\%)    & $95.8 \pm 0.0$ & $81.5 \pm 4.8$                           & $[+10.4, +18.2]$ \\
Metric-conflict recall (\%) & $0.0 \pm 0.0$ & $100.0 \pm 0.0$                         & --- \\
\bottomrule
\multicolumn{4}{l}{\textsuperscript{$\dagger$}Formal guarantee (Theorem~\ref{thm:safety}).}
\end{tabular}
\end{table*}

\subsection{Experiment 3: Full Degradation Table}
\begin{table*}[h]
\centering
\caption{Lineage-completeness degradation: leak rate vs.\ reporting completeness
$\rho$ (30 seeds each, mean $\pm$ population std dev).
At $\rho{=}1.0$ the guarantee holds by Theorem~\ref{thm:safety}.}
\label{tab:degradation}
\begin{tabular}{@{}lcccc@{}}
\toprule
\textbf{Schema} & $\rho=0.50$ & $\rho=0.75$ & $\rho=0.90$ & $\rho=1.00$ \\
\midrule
Synthetic (5-table) & $16.3 \pm 3.3\%$ & $0.0 \pm 0.0\%$ & $0.0 \pm 0.0\%$ & $0.0 \pm 0.0\%$ \\
TPC-H (8-table)     & $10.3 \pm 4.7\%$ & $0.9 \pm 1.1\%$ & $0.0 \pm 0.0\%$ & $0.0 \pm 0.0\%$ \\
\bottomrule
\end{tabular}
\end{table*}

\subsection{Experiment 4: Full Fuzzy Conflict-Detection Results}
\begin{table*}[h]
\centering
\caption{Extended fuzzy conflict-detection results (43 pairs:
15 TC + 15 TV + 8 LO + 5 CS). Bootstrap 95\% CI on $F_1$ shown in
square brackets. LO column shows per-category TP/total detection.}
\label{tab:fuzzy}
\begin{tabular}{@{}lccccc@{}}
\toprule
\textbf{Detector} & \textbf{Prec.} & \textbf{Recall} & \textbf{F\textsubscript{1}}
& \textbf{95\% CI ($F_1$)} & \textbf{LO correct} \\
\midrule
D1 --- Exact hash             & 0.651 & 1.000 & 0.789 & [0.677, 0.883] & 8/8\,\checkmark \\
D2 --- Jaccard ($\tau{=}0.70$)    & 0.676 & 0.821 & 0.742 & [0.604, 0.853] & 3/8 \\
D3 --- Column-graph           & \textbf{1.000} & 0.714 & 0.833 & [0.703, 0.933] & 0/8\,\ding{55} \\
\bottomrule
\end{tabular}
\end{table*}

\subsection{Experiment 5: Complete Runtime Measurements}
\begin{table*}[h]
\centering
\caption{Empirical runtime (mean $\pm$ std dev over 2,000 iterations).
$k{=}2$--$3$ lineage steps, $c{=}4$--$7$ columns per step.}
\label{tab:runtime}
\begin{tabular}{@{}llll@{}}
\toprule
\textbf{Operation} & \textbf{$n$} & \textbf{Mean ($\mu$s)} & \textbf{Std dev ($\mu$s)} \\
\midrule
Gate-check predicate (set-difference) & any  & $0.09$ & $0.03$ \\
Sensitivity-tag derivation ($O(kc)$)  & ---  & $0.35$ & $0.13$ \\
Definition-hash computation ($O(kc)$) & ---  & $1.20$ & $0.46$ \\
\midrule
Retrieve, best-case (first hit safe)  & 1    & $0.66$ & $0.23$ \\
Retrieve, best-case (first hit safe)  & 50   & $0.64$ & $0.22$ \\
Retrieve, worst-case (all blocked)    & 1    & $0.66$ & $0.12$ \\
Retrieve, worst-case (all blocked)    & 5    & $1.74$ & $0.22$ \\
Retrieve, worst-case (all blocked)    & 10   & $3.07$ & $0.39$ \\
Retrieve, worst-case (all blocked)    & 20   & $5.72$ & $0.45$ \\
Retrieve, worst-case (all blocked)    & 50   & $13.82$ & $0.76$ \\
\midrule
Write, no-conflict (full scan, $O(n)$)  & 1  & $0.28$ & $0.03$ \\
Write, no-conflict (full scan, $O(n)$)  & 50 & $1.63$ & $0.20$ \\
Write, conflict detected (early exit)   & 1  & $2.37$ & $1.12$ \\
Write, conflict detected (early exit)   & 50 & $2.33$ & $0.61$ \\
\bottomrule
\end{tabular}
\end{table*}

\subsection{Experiment 6: All Real-Agent Round-Trips}
\begin{table*}[h]
\centering
\caption{Experiment~6: all 9 agent round-trips. ``REUSE'' = served from
         cached AMU; ``WRITE'' = fresh SQL executed and AMU written;
         ``BLOCK'' = gate blocked at least one candidate before
         falling back to a fresh computation.}
\label{tab:exp6}
\small
\begin{tabular}{@{}p{0.03\linewidth}p{0.11\linewidth}p{0.24\linewidth}p{0.09\linewidth}p{0.1\linewidth}p{0.1\linewidth}p{0.1\linewidth}@{}}
\toprule
Label & Dept & Metric & Action & Blocked & Conflict & Leaked \\
\midrule
A1 & Finance    & revenue\_by\_category & WRITE & No & No & No \\
A2 & Sales      & revenue\_by\_category & REUSE & No & No & No \\
A3 & Operations & revenue\_by\_category & WRITE & \textbf{Yes} & \textbf{Yes} & No \\
B1 & HR         & employee\_directory   & WRITE & No & No & No \\
B2 & Finance    & employee\_directory   & REUSE & No & No & No \\
B3 & Sales      & employee\_directory   & WRITE & \textbf{Yes} & \textbf{Yes} & No \\
C1 & Finance    & orders\_by\_country   & WRITE & No & No & No \\
C2 & Operations & orders\_by\_country   & REUSE & No & No & No \\
C3 & Sales      & orders\_by\_country   & REUSE & No & No & No \\
\midrule
\multicolumn{4}{l}{\textit{Totals (9 round-trips)}} &
  2 blocked & 2 conflicts & \textbf{0 leaked} \\
\bottomrule
\end{tabular}
\end{table*}

% =============================================================================
\balance
\bibliographystyle{ieeetr}
\bibliography{references}
% =============================================================================
\end{document}
